\section{结论}

基于统一患者级 5-fold 公平比较协议，本文完成了 PTMC 术前 CentralLNM 预测证据链的系统重组。原始临床变量构成最稳定的参照基线；传统影像组学可提供补充信息，但未显示稳定超越；多模态深度学习在采用患者级 Attention-MIL、临床主导门控融合与先验约束后，获得了当前最具稳定性的深度模型结果，并相对于原始临床提供了有限但可重复的增量价值；多源融合则提示 ROI radiomics 仍具有进一步挖掘的空间。

因此，面向临床期刊的最终主线应当强调：在当前数据条件下，CentralLNM 的最佳建模思路不是抛开临床信息追求更复杂模型，而是在临床主路径基础上，以结构化、可解释且具有临床一致性的方式吸收影像信息。下一阶段若要进一步提升文章档次，重点不应停留在普通拼接融合，而应转向层级约束建模、患者级弱监督可解释 MIL 与个体化证据融合，并通过外部验证、增量统计检验和临床净获益分析证明其真实价值。
